\documentclass[10pt,a4paper]{article}

% ============================================================
%  Ordre d'inclusion : ne pas modifier
% ============================================================
% ============================================================
%  config.tex — Personnalisation du cours
%  Modifiez uniquement ce fichier pour adapter le template.
% ============================================================

% ------------------------------------------------------------
% Page de titre
% ------------------------------------------------------------
\newcommand{\coursTitre}{TITRE DE VOTRE COURS}
\newcommand{\coursSousTitre}{Sous-titre de votre cours}
\newcommand{\coursBandeau}{SUPPORT DE COURS · ÉDITION 2026}
\newcommand{\coursAuteur}{Prénom Nom}
\newcommand{\coursVersion}{Édition 2026}

% ------------------------------------------------------------
% En-tête de toutes les pages
% Gardez ce texte court pour éviter les débordements.
% ------------------------------------------------------------
\newcommand{\coursEntete}{Mon cours — Support de formation}

% ------------------------------------------------------------
% Bandeau affiché en haut des pages de chapitre
% Gardez ce texte court (une ligne).
% ------------------------------------------------------------
\newcommand{\coursBandeauChapitre}{SUPPORT DE COURS}

% ------------------------------------------------------------
% Taille du code (blocs de code + code inline)
% Valeurs possibles : \footnotesize, \small, \normalsize, \large
% ------------------------------------------------------------
\newcommand{\coursCodeSize}{\normalsize}

% ============================================================
%  preamble.tex — Moteur de style du cours
%
%  Ne modifiez pas ce fichier pour personnaliser votre cours :
%  utilisez config.tex.
%
%  Compilation : lualatex main.tex (2 passes)
% ============================================================

% ------------------------------------------------------------
% Mise en page et polices
% ------------------------------------------------------------
\usepackage[a4paper,margin=2cm,headheight=19pt]{geometry}
\usepackage{fontspec}
\setmainfont{Liberation Sans}
\setsansfont{Liberation Sans}
\setmonofont{JetBrains Mono}[
  Path=fonts/,
  UprightFont=JetBrainsMono-Regular.ttf,
  BoldFont=JetBrainsMono-Bold.ttf,
  ItalicFont=JetBrainsMono-Italic.ttf,
  BoldItalicFont=JetBrainsMono-BoldItalic.ttf,
  Scale=0.95
]

\usepackage[french]{babel}
\usepackage{microtype}
\usepackage{amsmath,amssymb}
\usepackage{graphicx}
\usepackage{float}
\usepackage{xcolor}
\usepackage{tikz}
\usetikzlibrary{arrows.meta,positioning,shapes.geometric,calc,fit,backgrounds}
\usepackage[most]{tcolorbox}
\tcbuselibrary{listings,skins,breakable}
\usepackage{listings}
\usepackage{enumitem}
\usepackage{booktabs}
\usepackage{tabularx}
\usepackage{array}
\usepackage{fancyhdr}
\usepackage{lastpage}
\usepackage{caption}
\usepackage{hyperref}

% ------------------------------------------------------------
% Couleurs
% ------------------------------------------------------------
\definecolor{deepblue}{HTML}{0F2A4A}
\definecolor{teal}{HTML}{0F766E}
\definecolor{warm}{HTML}{C2410C}
\definecolor{codebg}{HTML}{F6F8FA}
\definecolor{codegray}{HTML}{6B7280}
\definecolor{softblue}{HTML}{EAF2F8}
\definecolor{softgray}{HTML}{F4F6F8}
\definecolor{softteal}{HTML}{E6F4F1}

\hypersetup{
  colorlinks=true,
  linkcolor=deepblue,
  urlcolor=teal,
  citecolor=teal,
  pdftitle={\coursTitre\ — \coursSousTitre},
  pdfauthor={\coursAuteur},
  pdfsubject={\coursSousTitre},
  pdfkeywords={cours, support de cours, formation}
}

% ------------------------------------------------------------
% Typographie générale
% ------------------------------------------------------------
\setlength{\parindent}{0pt}
\setlength{\parskip}{0.45em}
\setlength{\emergencystretch}{2em}
\raggedbottom
\renewcommand{\arraystretch}{1.2}
\setlist[itemize]{leftmargin=1.4em,itemsep=2pt,topsep=3pt}
\setlist[enumerate]{leftmargin=1.7em,itemsep=2pt,topsep=3pt}

\captionsetup{font={small,sf},labelfont={bf,color=deepblue}}
\renewcommand{\lstlistingname}{Extrait de code}
\renewcommand{\lstlistlistingname}{Liste des extraits de code}

% ------------------------------------------------------------
% En-têtes et pieds de page
% ------------------------------------------------------------
\pagestyle{fancy}
\fancyhf{}
\fancyhead[L]{\sffamily\footnotesize\color{codegray}\coursEntete}
\fancyhead[R]{\sffamily\footnotesize\color{codegray}\leftmark}
\fancyfoot[C]{\sffamily\footnotesize\color{codegray}\thepage\ / \pageref{LastPage}}
\renewcommand{\headrulewidth}{0.4pt}
\renewcommand{\headrule}{\hbox to\headwidth{\color{teal!30}\leaders\hrule height \headrulewidth\hfill}}

% ------------------------------------------------------------
% Chapitres et sections personnalisés
% ------------------------------------------------------------
\newcounter{chapter}
\renewcommand{\thechapter}{\arabic{chapter}}

% \chapitre{numéro}{Titre}
% Le numéro est affiché. Il doit suivre l'ordre des chapitres.
\newcommand{\chapitre}[2]{%
  \clearpage
  \setcounter{chapter}{#1}%
  \phantomsection
  \pdfbookmark[1]{Chapitre #1\ — #2}{chapter#1}%
  \addcontentsline{toc}{section}{\protect\textbf{Chapitre #1\ — #2}}%
  \markboth{Chapitre #1\ — #2}{}%
  \thispagestyle{plain}%
  \begin{tikzpicture}[remember picture,overlay]
    \fill[deepblue] (current page.north west) rectangle
      ([yshift=-2.7cm]current page.north east);
    \fill[teal] ([yshift=-2.62cm]current page.north west) rectangle
      ([yshift=-2.7cm]current page.north east);
    \node[anchor=west,text=white,font=\sffamily\bfseries\fontsize{30}{34}\selectfont]
      at ([xshift=2cm,yshift=-1.15cm]current page.north west) {Chapitre #1};
    \node[anchor=west,text=teal!50,font=\sffamily\bfseries\fontsize{10}{12}\selectfont]
      at ([xshift=2cm,yshift=-2.05cm]current page.north west)
      {\addfontfeatures{LetterSpace=12}\coursBandeauChapitre};
  \end{tikzpicture}%
  \par\vspace*{0.7cm}
  {\sffamily\bfseries\color{deepblue}\fontsize{22}{26}\selectfont #2\par}
  \vspace{0.25cm}
  {\color{teal}\rule{\linewidth}{1.2pt}}
  \vspace{0.65cm}
}

% \soussection{numéro}{Titre}
\newcommand{\soussection}[2]{%
  \phantomsection
  \pdfbookmark[2]{#1 — #2}{sec:#1}%
  \addcontentsline{toc}{subsection}{\protect\numberline{#1}#2}%
  \markright{#1 — #2}%
  {\sffamily\bfseries\color{teal}\large #1\quad #2\par}
  \vspace{0.25cm}
}

% \soussoussection{numéro}{Titre}
\newcommand{\soussoussection}[2]{%
  \phantomsection
  \pdfbookmark[3]{#1 — #2}{subsec:#1}%
  \addcontentsline{toc}{subsubsection}{\protect\numberline{#1}#2}%
  {\sffamily\bfseries\color{deepblue}\normalsize #1\quad #2\par}
  \vspace{0.15cm}
}

% ------------------------------------------------------------
% Encadrés graphiques
% Tous s'utilisent de la même façon :
%   \begin{nom}{} ... \end{nom}
% ------------------------------------------------------------
\newtcolorbox{objectifs}[1][]{
  enhanced,breakable,
  title={\sffamily\bfseries Objectifs du chapitre},
  colback=deepblue!2,colframe=deepblue!35,
  colbacktitle=deepblue!8,coltitle=deepblue,
  fonttitle=\small,boxrule=0.5pt,arc=1.5mm,
  left=8pt,right=8pt,top=7pt,bottom=7pt,#1
}

\newtcolorbox{definition}[1][]{
  enhanced,breakable,
  title={\sffamily\bfseries Définition},
  colback=white,colframe=deepblue!30,
  colbacktitle=deepblue!6,coltitle=deepblue,
  fonttitle=\small,boxrule=0.5pt,arc=1.5mm,
  left=8pt,right=8pt,top=7pt,bottom=7pt,#1
}

\newtcolorbox{retenir}[1][]{
  enhanced,breakable,
  title={\sffamily\bfseries À retenir},
  colback=white,colframe=teal!45,
  colbacktitle=teal!8,coltitle=teal!80!black,
  fonttitle=\small,boxrule=0.5pt,arc=1.5mm,
  left=8pt,right=8pt,top=7pt,bottom=7pt,#1
}

\newtcolorbox{attention}[1][]{
  enhanced,breakable,
  title={\sffamily\bfseries Attention},
  colback=red!3,colframe=red!45!black,
  colbacktitle=red!8,coltitle=red!60!black,
  fonttitle=\small,boxrule=0.5pt,arc=1.5mm,
  left=8pt,right=8pt,top=7pt,bottom=7pt,#1
}

\newtcolorbox{piege}[1][]{
  enhanced,breakable,
  title={\sffamily\bfseries Piège classique},
  colback=warm!3,colframe=warm!45,
  colbacktitle=warm!9,coltitle=warm!80!black,
  fonttitle=\small,boxrule=0.5pt,arc=1.5mm,
  left=8pt,right=8pt,top=7pt,bottom=7pt,#1
}

\newtcolorbox{conseil}[1][]{
  enhanced,breakable,
  title={\sffamily\bfseries Conseil},
  colback=teal!3,colframe=teal!45,
  colbacktitle=teal!10,coltitle=teal!80!black,
  fonttitle=\small,boxrule=0.5pt,arc=1.5mm,
  left=8pt,right=8pt,top=7pt,bottom=7pt,#1
}

\newtcolorbox{choix}[1][]{
  enhanced,breakable,
  title={\sffamily\bfseries Choix recommandé},
  colback=softblue,colframe=deepblue!40,
  colbacktitle=deepblue!8,coltitle=deepblue,
  fonttitle=\small,boxrule=0.5pt,arc=1.5mm,
  left=8pt,right=8pt,top=7pt,bottom=7pt,#1
}

\newtcolorbox{atelier}[1][]{
  enhanced,breakable,
  title={\sffamily\bfseries Atelier},
  colback=softgray,colframe=codegray!45,
  colbacktitle=codegray!12,coltitle=black!75,
  fonttitle=\small,boxrule=0.5pt,arc=1.5mm,
  left=8pt,right=8pt,top=7pt,bottom=7pt,#1
}

\newtcolorbox{resume}[1][]{
  enhanced,breakable,
  title={\sffamily\bfseries En résumé},
  colback=softteal,colframe=teal!25,
  colbacktitle=teal!8,coltitle=teal!80!black,
  fonttitle=\small,boxrule=0.4pt,arc=1.2mm,
  left=8pt,right=8pt,top=6pt,bottom=6pt,#1
}

% ------------------------------------------------------------
% Langages et styles de code
% ------------------------------------------------------------
\lstdefinelanguage{Java25}{
  language=Java,
  morekeywords={
    record, sealed, permits, non-sealed, yield, var, when
  },
  morekeywords=[2]{
    @SpringBootApplication, @RestController, @Controller, @Service, @Repository,
    @Component, @Configuration, @Bean, @Autowired, @Value, @ConfigurationProperties,
    @GetMapping, @PostMapping, @PutMapping, @DeleteMapping, @PatchMapping,
    @RequestBody, @PathVariable, @RequestParam, @ResponseStatus,
    @RestControllerAdvice, @ControllerAdvice, @ExceptionHandler,
    @CrossOrigin, @PostConstruct, @PreDestroy, @Valid, @NotNull, @NotBlank,
    @Size, @Min, @Max, @Tag, @Operation, @ApiResponse, @Parameter, @Schema
  }
}

\lstdefinelanguage{GradleGroovy}{
  keywords={
    plugins, id, version, repositories, mavenCentral, dependencies,
    implementation, testImplementation, runtimeOnly, compileOnly,
    java, toolchain, languageVersion, JavaLanguageVersion,
    group, sourceCompatibility, tasks, withType, named, test
  }
}

\lstdefinelanguage{PropertiesLang}{
  keywords={},
  sensitive=false
}

% Style commun à tous les blocs de code
\lstdefinestyle{coursbase}{
  backgroundcolor=\color{codebg},
  basicstyle=\ttfamily\bfseries\coursCodeSize,
  keywordstyle=\color{deepblue}\bfseries,
  keywordstyle=[2]\color{teal}\bfseries,
  commentstyle=\color{teal!70!black}\itshape,
  stringstyle=\color{warm!80!black},
  numbers=none,
  frame=single,rulecolor=\color{deepblue!20},
  breaklines=true,breakatwhitespace=false,
  postbreak=\mbox{\textcolor{teal}{$\hookrightarrow$}\space},
  showstringspaces=false,tabsize=2,columns=fullflexible,
  keepspaces=true,upquote=true,
  captionpos=b,xleftmargin=0.5em,framexleftmargin=0.3em,
  lineskip=-0.6pt,
  aboveskip=4pt,belowskip=4pt
}

\lstdefinestyle{java}{style=coursbase,language=Java25}
\lstdefinestyle{gradle}{style=coursbase,language=GradleGroovy}
\lstdefinestyle{properties}{style=coursbase,language=PropertiesLang}
\lstdefinestyle{bash}{style=coursbase,language=bash}
\lstdefinestyle{json}{style=coursbase}
\lstdefinestyle{sql}{style=coursbase,language=SQL}
\lstdefinestyle{python}{style=coursbase,language=Python}
\lstdefinestyle{html}{style=coursbase,language=HTML}
\lstdefinestyle{xml}{style=coursbase,language=XML}
\lstdefinestyle{c}{style=coursbase,language=C}
\lstdefinestyle{cpp}{style=coursbase,language=C++}
\lstdefinestyle{text}{style=coursbase}

% ------------------------------------------------------------
% Raccourcis de mise en forme
% ------------------------------------------------------------
% Code court dans le texte : \code{nom_de_variable}
\DeclareRobustCommand{\code}[1]{%
  \texorpdfstring{\colorbox{codebg}{\texttt{\coursCodeSize\bfseries\color{deepblue!90}\detokenize{#1}}}}{\detokenize{#1}}%
}

% Terme important : \motcle{concept}
\newcommand{\motcle}[1]{\textsf{\color{teal!80!black}#1}}

% Alias de compatibilité avec l'ancien template
\newcommand{\spring}[1]{\motcle{#1}}

% ------------------------------------------------------------
% Table des matières : largeur des numéros
% ------------------------------------------------------------
\makeatletter
\renewcommand*\l@section{\@dottedtocline{1}{0.5em}{2.8em}}
\renewcommand*\l@subsection{\@dottedtocline{2}{2.0em}{3.4em}}
\renewcommand*\l@subsubsection{\@dottedtocline{3}{4.2em}{4.2em}}
\makeatother


\begin{document}

% Pages liminaires : numérotation romaine, sans en-tête ni pied de page
\pagenumbering{roman}
\pagestyle{empty}

% ============================================================
% Page de titre
% Les textes sont définis dans config.tex
% ============================================================
\begin{titlepage}
\begin{tikzpicture}[remember picture,overlay]
  \fill[deepblue] (current page.north west) rectangle (current page.south east);
  \fill[teal!85!black] (current page.north west) rectangle
    ([yshift=-0.35cm]current page.north east);
  \fill[teal!25] ([yshift=3.2cm]current page.south west) rectangle
    ([yshift=3.35cm]current page.south east);

  \node[anchor=west,text=teal!35,font=\sffamily\bfseries\fontsize{13}{16}\selectfont]
    at ([xshift=2.2cm,yshift=-4.0cm]current page.north west)
    {\coursBandeau};

  \node[anchor=north west,text=white,font=\sffamily\bfseries\fontsize{30}{34}\selectfont,
        text width=16.2cm,align=left] (titre)
    at ([xshift=2.2cm,yshift=-4.8cm]current page.north west)
    {\coursTitre};

  \node[anchor=north west,text=teal!30,font=\sffamily\bfseries\fontsize{18}{22}\selectfont,
        text width=16.2cm,align=left]
    at ([yshift=-0.4cm]titre.south west)
    {\coursSousTitre};

  \node[anchor=west,text=white,font=\sffamily\fontsize{12}{15}\selectfont]
    at ([xshift=2.2cm,yshift=2.55cm]current page.south west)
    {\coursAuteur};

  \node[anchor=east,text=teal!35,font=\sffamily\bfseries\fontsize{11}{14}\selectfont]
    at ([xshift=-2.2cm,yshift=2.55cm]current page.south east)
    {\coursVersion};
\end{tikzpicture}
\end{titlepage}

% ============================================================
% Table des matières
% ============================================================
\clearpage
{\small
\tableofcontents
}
\clearpage

% ============================================================
% Chapitres
% Ajoutez vos chapitres ici, dans l'ordre.
% ============================================================
\pagenumbering{arabic}
\pagestyle{fancy}
\setcounter{page}{1}

% ============================================================
%  CHAPITRE 1 — À COMPLÉTER
%
%  Remplacez les textes entre [crochets] par votre contenu.
%  Pour ajouter un chapitre, copiez plutôt :
%      chapitres/modele-chapitre.tex
% ============================================================

\chapitre{1}{Titre du chapitre 1}

\begin{objectifs}
\begin{itemize}
  \item Comprendre [premier objectif du chapitre].
  \item Savoir [deuxième objectif du chapitre].
  \item Être capable de [troisième objectif du chapitre].
\end{itemize}
\end{objectifs}

\soussection{1.1}{Titre de la section 1.1}

Rédigez ici l'introduction, la description ou l'explication de votre
première section. Un paragraphe normal s'écrit directement dans le
fichier, sans commande particulière.

\soussoussection{1.1.1}{Titre de la sous-section 1.1.1}

Rédigez ici le premier point détaillé. Utilisez \code{un_nom_de_variable}
pour les éléments techniques et \motcle{un terme clé} pour attirer
l'attention sur un concept important.

\soussoussection{1.1.2}{Titre de la sous-section 1.1.2}

Ajoutez ici un deuxième point. Vous pouvez insérer un encadré, un bloc
de code, un tableau ou une image. Tous les exemples sont disponibles
dans le fichier \code{chapitres/demo.tex}.

\soussection{1.2}{Titre de la section 1.2}

Rédigez ici la deuxième partie du chapitre.

\begin{resume}
Résumez en quelques lignes les points essentiels du chapitre 1.
\end{resume}

% ============================================================
%  CHAPITRE 2 — À COMPLÉTER
% ============================================================

\chapitre{2}{Titre du chapitre 2}

\begin{objectifs}
\begin{itemize}
  \item Comprendre [premier objectif du chapitre].
  \item Savoir [deuxième objectif du chapitre].
\end{itemize}
\end{objectifs}

\soussection{2.1}{Titre de la section 2.1}

Rédigez ici le contenu du chapitre 2.

\soussoussection{2.1.1}{Titre de la sous-section 2.1.1}

Rédigez ici le contenu détaillé.

\soussection{2.2}{Titre de la section 2.2}

Rédigez ici la suite du chapitre.

\begin{resume}
Résumez en quelques lignes les points essentiels du chapitre 2.
\end{resume}


% Ajoutez vos nouveaux chapitres ici, par exemple :
% \input{chapitres/ch3}

% Décommentez la ligne suivante pour afficher le chapitre
% de démonstration (toutes les fonctionnalités du template).
% % ============================================================
%  CHAPITRE DE DÉMONSTRATION
%
%  Ce fichier n'est pas inclus par défaut.
%  Pour l'afficher : décommentez % ============================================================
%  CHAPITRE DE DÉMONSTRATION
%
%  Ce fichier n'est pas inclus par défaut.
%  Pour l'afficher : décommentez % ============================================================
%  CHAPITRE DE DÉMONSTRATION
%
%  Ce fichier n'est pas inclus par défaut.
%  Pour l'afficher : décommentez \input{chapitres/demo} dans main.tex.
%  Il montre tous les éléments disponibles dans le template.
% ============================================================

\chapitre{3}{Démonstration : toutes les fonctionnalités}

\begin{objectifs}
\begin{itemize}
  \item Voir un exemple de chaque élément disponible.
  \item Savoir quel environnement utiliser pour chaque besoin.
\end{itemize}
\end{objectifs}

\soussection{3.1}{Texte, listes et mise en forme}

Un paragraphe normal s'écrit directement. On peut mettre un mot en
\textbf{gras}, un autre en \textit{italique}, ou afficher du
\code{code inline}. Le terme \motcle{important} est coloré.

\begin{itemize}
  \item Premier point.
  \item Deuxième point avec une liste imbriquée :
  \begin{enumerate}
    \item Sous-point un.
    \item Sous-point deux.
  \end{enumerate}
\end{itemize}

\begin{definition}
Une \textbf{définition} présente un concept nouveau de manière rigoureuse
et précise.
\end{definition}

\begin{retenir}
Ce bloc sert à résumer ce qu'il faut absolument retenir.
\end{retenir}

\begin{attention}
Ce bloc signale un point dangereux ou une erreur à ne pas commettre.
\end{attention}

\begin{piege}
Ce bloc décrit un piège classique et la manière de l'éviter.
\end{piege}

\begin{conseil}
Ce bloc donne une bonne pratique ou une astuce.
\end{conseil}

\begin{choix}
Ce bloc justifie un choix technique recommandé.
\end{choix}

\begin{atelier}
Ce bloc propose une activité pratique ou un exercice guidé.
\end{atelier}

\soussection{3.2}{Blocs de code}

Voici un exemple Java :

\begin{lstlisting}[style=java,caption={Exemple de contrôleur Java}]
@RestController
@RequestMapping("/api/articles")
public class ArticleController {

    @GetMapping("/{id}")
    public ResponseEntity<Article> findById(@PathVariable Long id) {
        return ResponseEntity.ok(service.findById(id));
    }
}
\end{lstlisting}

Un exemple Python :

\begin{lstlisting}[style=python,caption={Exemple Python}]
def salutation(nom: str) -> str:
    return f"Bonjour, {nom} !"

print(salutation("Alice"))
\end{lstlisting}

Un exemple JSON :

\begin{lstlisting}[style=json,caption={Exemple de réponse JSON}]
{
  "id": 42,
  "titre": "Mon article",
  "publie": true
}
\end{lstlisting}

Une commande shell :

\begin{lstlisting}[style=bash,caption={Commandes de compilation}]
lualatex main.tex
lualatex main.tex
\end{lstlisting}

Un fichier de configuration :

\begin{lstlisting}[style=properties,caption={Fichier de configuration}]
serveur.port=8080
serveur.hote=localhost
mode.debug=true
\end{lstlisting}

\soussection{3.3}{Tableaux et mathématiques}

\begin{table}[htbp]
\centering
\small
\begin{tabularx}{\textwidth}{lX}
\toprule
\textbf{Élément} & \textbf{Description} \\
\midrule
\code{lstlisting} & Affiche un bloc de code. \\
\code{tcolorbox} & Affiche un encadré coloré. \\
\code{tabularx} & Affiche un tableau qui s'adapte à la largeur. \\
\bottomrule
\end{tabularx}
\caption{Exemple de tableau avec booktabs et tabularx}
\end{table}

Une formule mathématique :

\begin{equation*}
  f(x) = \sum_{i=1}^{n} x_i^2
\end{equation*}

Un lien : \href{https://example.com}{https://example.com}.

\soussection{3.4}{Conclusion}

\begin{resume}
Ce chapitre de démonstration montre tous les éléments de base du
template. Il n'est pas destiné à être conservé dans le cours final.
\end{resume}
 dans main.tex.
%  Il montre tous les éléments disponibles dans le template.
% ============================================================

\chapitre{3}{Démonstration : toutes les fonctionnalités}

\begin{objectifs}
\begin{itemize}
  \item Voir un exemple de chaque élément disponible.
  \item Savoir quel environnement utiliser pour chaque besoin.
\end{itemize}
\end{objectifs}

\soussection{3.1}{Texte, listes et mise en forme}

Un paragraphe normal s'écrit directement. On peut mettre un mot en
\textbf{gras}, un autre en \textit{italique}, ou afficher du
\code{code inline}. Le terme \motcle{important} est coloré.

\begin{itemize}
  \item Premier point.
  \item Deuxième point avec une liste imbriquée :
  \begin{enumerate}
    \item Sous-point un.
    \item Sous-point deux.
  \end{enumerate}
\end{itemize}

\begin{definition}
Une \textbf{définition} présente un concept nouveau de manière rigoureuse
et précise.
\end{definition}

\begin{retenir}
Ce bloc sert à résumer ce qu'il faut absolument retenir.
\end{retenir}

\begin{attention}
Ce bloc signale un point dangereux ou une erreur à ne pas commettre.
\end{attention}

\begin{piege}
Ce bloc décrit un piège classique et la manière de l'éviter.
\end{piege}

\begin{conseil}
Ce bloc donne une bonne pratique ou une astuce.
\end{conseil}

\begin{choix}
Ce bloc justifie un choix technique recommandé.
\end{choix}

\begin{atelier}
Ce bloc propose une activité pratique ou un exercice guidé.
\end{atelier}

\soussection{3.2}{Blocs de code}

Voici un exemple Java :

\begin{lstlisting}[style=java,caption={Exemple de contrôleur Java}]
@RestController
@RequestMapping("/api/articles")
public class ArticleController {

    @GetMapping("/{id}")
    public ResponseEntity<Article> findById(@PathVariable Long id) {
        return ResponseEntity.ok(service.findById(id));
    }
}
\end{lstlisting}

Un exemple Python :

\begin{lstlisting}[style=python,caption={Exemple Python}]
def salutation(nom: str) -> str:
    return f"Bonjour, {nom} !"

print(salutation("Alice"))
\end{lstlisting}

Un exemple JSON :

\begin{lstlisting}[style=json,caption={Exemple de réponse JSON}]
{
  "id": 42,
  "titre": "Mon article",
  "publie": true
}
\end{lstlisting}

Une commande shell :

\begin{lstlisting}[style=bash,caption={Commandes de compilation}]
lualatex main.tex
lualatex main.tex
\end{lstlisting}

Un fichier de configuration :

\begin{lstlisting}[style=properties,caption={Fichier de configuration}]
serveur.port=8080
serveur.hote=localhost
mode.debug=true
\end{lstlisting}

\soussection{3.3}{Tableaux et mathématiques}

\begin{table}[htbp]
\centering
\small
\begin{tabularx}{\textwidth}{lX}
\toprule
\textbf{Élément} & \textbf{Description} \\
\midrule
\code{lstlisting} & Affiche un bloc de code. \\
\code{tcolorbox} & Affiche un encadré coloré. \\
\code{tabularx} & Affiche un tableau qui s'adapte à la largeur. \\
\bottomrule
\end{tabularx}
\caption{Exemple de tableau avec booktabs et tabularx}
\end{table}

Une formule mathématique :

\begin{equation*}
  f(x) = \sum_{i=1}^{n} x_i^2
\end{equation*}

Un lien : \href{https://example.com}{https://example.com}.

\soussection{3.4}{Conclusion}

\begin{resume}
Ce chapitre de démonstration montre tous les éléments de base du
template. Il n'est pas destiné à être conservé dans le cours final.
\end{resume}
 dans main.tex.
%  Il montre tous les éléments disponibles dans le template.
% ============================================================

\chapitre{3}{Démonstration : toutes les fonctionnalités}

\begin{objectifs}
\begin{itemize}
  \item Voir un exemple de chaque élément disponible.
  \item Savoir quel environnement utiliser pour chaque besoin.
\end{itemize}
\end{objectifs}

\soussection{3.1}{Texte, listes et mise en forme}

Un paragraphe normal s'écrit directement. On peut mettre un mot en
\textbf{gras}, un autre en \textit{italique}, ou afficher du
\code{code inline}. Le terme \motcle{important} est coloré.

\begin{itemize}
  \item Premier point.
  \item Deuxième point avec une liste imbriquée :
  \begin{enumerate}
    \item Sous-point un.
    \item Sous-point deux.
  \end{enumerate}
\end{itemize}

\begin{definition}
Une \textbf{définition} présente un concept nouveau de manière rigoureuse
et précise.
\end{definition}

\begin{retenir}
Ce bloc sert à résumer ce qu'il faut absolument retenir.
\end{retenir}

\begin{attention}
Ce bloc signale un point dangereux ou une erreur à ne pas commettre.
\end{attention}

\begin{piege}
Ce bloc décrit un piège classique et la manière de l'éviter.
\end{piege}

\begin{conseil}
Ce bloc donne une bonne pratique ou une astuce.
\end{conseil}

\begin{choix}
Ce bloc justifie un choix technique recommandé.
\end{choix}

\begin{atelier}
Ce bloc propose une activité pratique ou un exercice guidé.
\end{atelier}

\soussection{3.2}{Blocs de code}

Voici un exemple Java :

\begin{lstlisting}[style=java,caption={Exemple de contrôleur Java}]
@RestController
@RequestMapping("/api/articles")
public class ArticleController {

    @GetMapping("/{id}")
    public ResponseEntity<Article> findById(@PathVariable Long id) {
        return ResponseEntity.ok(service.findById(id));
    }
}
\end{lstlisting}

Un exemple Python :

\begin{lstlisting}[style=python,caption={Exemple Python}]
def salutation(nom: str) -> str:
    return f"Bonjour, {nom} !"

print(salutation("Alice"))
\end{lstlisting}

Un exemple JSON :

\begin{lstlisting}[style=json,caption={Exemple de réponse JSON}]
{
  "id": 42,
  "titre": "Mon article",
  "publie": true
}
\end{lstlisting}

Une commande shell :

\begin{lstlisting}[style=bash,caption={Commandes de compilation}]
lualatex main.tex
lualatex main.tex
\end{lstlisting}

Un fichier de configuration :

\begin{lstlisting}[style=properties,caption={Fichier de configuration}]
serveur.port=8080
serveur.hote=localhost
mode.debug=true
\end{lstlisting}

\soussection{3.3}{Tableaux et mathématiques}

\begin{table}[htbp]
\centering
\small
\begin{tabularx}{\textwidth}{lX}
\toprule
\textbf{Élément} & \textbf{Description} \\
\midrule
\code{lstlisting} & Affiche un bloc de code. \\
\code{tcolorbox} & Affiche un encadré coloré. \\
\code{tabularx} & Affiche un tableau qui s'adapte à la largeur. \\
\bottomrule
\end{tabularx}
\caption{Exemple de tableau avec booktabs et tabularx}
\end{table}

Une formule mathématique :

\begin{equation*}
  f(x) = \sum_{i=1}^{n} x_i^2
\end{equation*}

Un lien : \href{https://example.com}{https://example.com}.

\soussection{3.4}{Conclusion}

\begin{resume}
Ce chapitre de démonstration montre tous les éléments de base du
template. Il n'est pas destiné à être conservé dans le cours final.
\end{resume}


\end{document}
